\section{Objetivos y Preguntas de Investigación}

\subsection{Preguntas de Investigación}
\textbf{Pregunta General:}
¿De qué manera ha evolucionado la relación entre los ingresos de los hogares y el costo de vida en Chile, Brasil y Argentina durante el período 2015--2025, y cómo impacta esta dinámica según estratos socioeconómicos de ingreso?

\textbf{Preguntas Específicas:}
\begin{enumerate}
    \item ¿Cómo se comportan los ingresos de los hogares y el valor de la Canasta Básica Alimentaria (CBA) y Total (CBT) en cada uno de los tres países durante el decenio estudiado 2015--2025?
    \item ¿Existen diferencias significativas en la pérdida o ganancia del poder adquisitivo real entre los distintos quintiles de ingreso entre Chile, Brasil y Argentina?
    \item ¿Qué influencia ejercen las diferencias en las políticas institucionales de transferencia monetaria y las variaciones inflacionarias de cada país sobre el poder de compra real?
\end{enumerate}

\subsection{Objetivos}
\textbf{Objetivo General:}
Analizar la relación entre la evolución de los ingresos de los hogares y el costo de vida en Chile, Brasil y Argentina durante el decenio 2015--2025, evaluando las variaciones en el poder adquisitivo y las heterogeneidades socioeconómicas entre quintiles de ingreso.

\textbf{Objetivos Específicos:}
\begin{enumerate}
    \item Procesar y estandarizar los microdatos de las encuestas de hogares de Chile (\textit{CASEN}), Brasil (\textit{PNAD Contínua}) y Argentina (\textit{EPH}) para el período 2015--2025, equiparando las variables de ingresos y características sociodemográficas del hogar.
    \item Definir y calcular los indicadores de costo de vida, IPC y canastas básicas (CBA y CBT) para cada país dentro del marco temporal seleccionado.
    \item Calcular la variación del poder adquisitivo real por quintil de ingreso en las tres economías, identificando brechas y patrones frente a la inflación y shocks macroeconómicos.
    \item Comparar socioeconómicamente las trayectorias de Chile, Brasil y Argentina a la luz de sus marcos institucionales y de política de transferencias.
\end{enumerate}